\documentclass[11pt,a4paper]{article}

% ========================================
% PACKAGES
% ========================================
\usepackage[utf8]{inputenc}
\usepackage[T1]{fontenc}
\usepackage{lmodern}  % Better font support for larger sizes
\usepackage[english]{babel}  % Change to your language
\usepackage[margin=2.5cm]{geometry}

% Mathematics
\usepackage{amsmath, amssymb, amsthm}
\usepackage{mathtools}

% Graphics and colors
\usepackage{graphicx}
\usepackage{xcolor}
\usepackage{tikz}
\usetikzlibrary{positioning, arrows.meta, shapes.geometric, calc}
\usepackage{pgfplots}
\pgfplotsset{compat=1.18}
\usepackage[most]{tcolorbox}

% Lists and formatting
\usepackage{enumitem}
\usepackage{parskip}
\usepackage{fancyhdr}

% Code listings (if needed)
\usepackage{listings}
\usepackage{algorithm}
\usepackage{algorithmic}

% Hyperlinks
\usepackage{hyperref}
\hypersetup{
    colorlinks=true,
    linkcolor=blue,
    urlcolor=blue,
    citecolor=blue
}

% ========================================
% THEOREM ENVIRONMENTS
% ========================================
\theoremstyle{definition}
\newtheorem{definition}{Definition}[section]
\newtheorem{example}{Example}[section]
\newtheorem{exercise}{Exercise}[section]

\theoremstyle{plain}
\newtheorem{theorem}{Theorem}[section]
\newtheorem{lemma}[theorem]{Lemma}
\newtheorem{proposition}[theorem]{Proposition}
\newtheorem{corollary}[theorem]{Corollary}

\theoremstyle{remark}
\newtheorem*{remark}{Remark}
\newtheorem*{observation}{Observation}

% ========================================
% CUSTOM COMMANDS
% ========================================
\newcommand{\R}{\mathbb{R}}
\newcommand{\N}{\mathbb{N}}
\newcommand{\Z}{\mathbb{Z}}
\newcommand{\Q}{\mathbb{Q}}
\newcommand{\C}{\mathbb{C}}

% ========================================
% TABLE OF CONTENTS DEPTH
% ========================================
\setcounter{tocdepth}{2} % Show only parts, sections, and subsections

% ========================================
% HEADER AND FOOTER
% ========================================
\setlength{\headheight}{14pt}
\pagestyle{fancy}
\fancyhf{}
\lhead{\leftmark}
\rhead{PL \& RL}
\cfoot{\thepage}
\renewcommand{\sectionmark}[1]{\markboth{#1}{}}

% ========================================
% DOCUMENT INFORMATION
% ========================================
\title{\textbf{Planning \& Reinforcement Learning - Part 2}\\
\large Artificial Intelligence}
\author{Jacopo Parretti}
\date{I Semester 2025-2026}

% ========================================
% DOCUMENT
% ========================================
\begin{document}

% Custom title page
\begin{titlepage}
    \begin{center}
        \vspace*{2cm}

        % University/Course info at top
        {\large\textsc{Artificial Intelligence}}

        \vspace{0.4cm}
        \textcolor{blue!60!black}{\rule{0.5\textwidth}{1pt}}

        \vspace{2.5cm}

        % Main title with modern spacing
        {\fontsize{36}{44}\selectfont\bfseries Planning \&\\[0.5cm] Reinforcement Learning\\[0.5cm] Part 2}

        \vspace{1.5cm}

        % Subtitle with clean design
        \begin{tcolorbox}[
            colback=blue!3!white,
            colframe=blue!40!black,
            width=0.6\textwidth,
            arc=1mm,
            boxrule=0.8pt,
            left=15pt,
            right=15pt,
            top=10pt,
            bottom=10pt
        ]
            \centering
            {\large\textit{Lecture Notes}}\\[0.3cm]
            \textcolor{gray}{Master's Program in Artificial Intelligence}
        \end{tcolorbox}

        \vfill

        % Clean course details
        \begin{tcolorbox}[
            enhanced,
            colback=white,
            colframe=gray!30,
            width=0.5\textwidth,
            arc=0mm,
            boxrule=0.5pt,
            left=20pt,
            right=20pt,
            top=12pt,
            bottom=12pt,
            shadow={2mm}{-2mm}{0mm}{black!10}
        ]
            \centering
            {\small\textcolor{gray!70}{\textbf{ACADEMIC YEAR}}}\\[0.4cm]
            {\Large 2025--2026}
        \end{tcolorbox}

        \vspace{2cm}

        % Author
        {\LARGE\textsc{Jacopo Parretti}}

        \vspace{1.5cm}

    \end{center}
\end{titlepage}

\newpage
\tableofcontents
\newpage

\part{}

\section{}

State-Transition Systems

We want to model a complex environment - we need simplifying assumtions
classical planning:
- finite, static workd with just one actor
- no concurrent actions, no explicit time
- determinism, no uncertainty

sequence of states and actions <s_0, a_1, s_1, a_2, ...>

most real world environment violate these assumption but classical planning could be a useful approximation

State transition systems: \Epsilon = < S, A, \gamma, cost >
S and A, finite set of states and actions

plan: a sequence of actions \pi = <a_1, a_2, ..., a_n>

empty plan <>
concatenation of plans \pi \circ \pi' 

\gamma(s_0, \pi) = s_n means that after executing \pi starting from s_0 we end up in s_n 

run-plan(\Epsilon, \pi) 
while True do
    \dots


given a transition system, a planning problem is P = (\Epsilon, s_0, S_g) 
from s_0 to one of the states in S_g 

solution for P:
- a plan \pi such that \gamma(s_0, \pi) \in S_g 
- \pi is minimal if no subsequence of \pi is a solution
- pi is shortest if no solution pi', where |pi'| < |pi| 
- pi is optimal if cost(pi) = min{cost(pi') | pi' is a solution for P}


Representing a transition system \Epsilon 

(small)
if S and A are small enough ---> state/action unique name ---> for each s, a, \dots

(large)
cannot represent all states explicitly 


the more we are specific to the representation, the more we can reason 
high effectiveness: can use whatever works best for specific domain
low generality tho 

domain independent:


\dots


State-varaible representations 

example of robots, containers, locs 

r1 is a robot, but actually r1 is a symbol 
we consider it a robot tho 

symbols for objects/actions should interpret correctly the environment
 

rigid properties: properties of the states that don't change over time 

varying properties (fluents): represented with a state variable; a term to which we can assign a value

states s are functions that assigns values to state variables
one state is an assingment 
can be also a set of ground positive literals (= is a used as a predicate symbol)


action template/schema: a parametrized set of actions 


actions: \A = set of action templates 
assign each param a value in its range 











\end{document}
